\documentclass[11pt,a4paper]{article}

% --- Packages ---
\usepackage[utf8]{inputenc}
\usepackage[T1]{fontenc}
\usepackage{amsmath,amssymb,amsthm}
\usepackage{mathtools}
\usepackage{booktabs}
\usepackage{tabularx}
\usepackage{hyperref}
\usepackage{cleveref}
\usepackage{algorithm}
\usepackage{algorithmic}
\usepackage{graphicx}
\usepackage{xcolor}
\usepackage{listings}
\usepackage[margin=1in]{geometry}
\usepackage{enumitem}

% --- Theorem environments ---
\newtheorem{theorem}{Theorem}
\newtheorem{definition}{Definition}

% --- Listings style ---
\lstset{
  basicstyle=\ttfamily\small,
  keywordstyle=\color{blue},
  commentstyle=\color{gray},
  stringstyle=\color{red!60!black},
  breaklines=true,
  frame=single,
  columns=fullflexible,
}

% --- Title ---
\title{%
  \textbf{tqai: Efficient KV Cache Compression with Incremental Reconstruction,\\
  Fused Metal Kernels, Composable Pipeline Architecture,\\
  and Adaptive Bit Allocation for Local LLM and DiT Inference}%
}

\author{
  Patrick Bertsch \and Claude (Anthropic)
}

\date{April 2026 (v0.4 revision)}

\begin{document}

\maketitle

% ============================================================================
% ABSTRACT
% ============================================================================
\begin{abstract}
We present \texttt{tqai}, a comprehensive KV cache compression library for local inference of large language models (LLMs) and diffusion transformers (DiTs) on Apple Silicon. Building on the TurboQuant framework~\cite{zandieh2026turboquant}, \texttt{tqai} introduces five key contributions:
(1)~an \textbf{incremental cache reconstruction} strategy that reduces per-token dequantization from $O(n)$ to $O(1)$, delivering $2.3$--$3.6\times$ throughput improvement over naive full-reconstruction approaches;
(2)~\textbf{fused Metal GPU kernels} that eliminate Python dispatch overhead by combining $L_2$ normalization, orthogonal rotation, and codebook search into single GPU dispatches via MLX's \texttt{mx.fast.metal\_kernel} API;
(3)~\textbf{evolutionary codebook optimization} using CMA-ES and fuzzy C-means with temperature annealing, achieving $0.55\%$ MSE improvement over classical Lloyd-Max codebooks;
(4)~a \textbf{composable middleware pipeline (v0.4)} that turns the library from a single-algorithm compressor into a plugin framework where each new compression paper --- TurboQuant, QuantSparse, DiTFastAttn, BSA, Sheaf, Copresheaf, Fisher, SparseDiT --- ships as a single file without modifying the proven core;
and (5)~\textbf{production diffusers video pipeline support} including a VAE memory optimization that prevents the ${>}100$~GB peak that otherwise OOMs WAN~2.2 5B at 81-frame video on a 128~GB Mac, plus an MPS \texttt{float64} fix that unblocks LTX-2 on Apple Silicon.
We additionally ship a \textbf{chunked attention} implementation using the online softmax algorithm, a \textbf{Palm information-theoretic scorer} for adaptive bit allocation, and \textbf{CFG attention sharing} hooks.
Across 12 model variants ($0.5$B--$14$B parameters) including Gemma~4, \texttt{tqai} achieves \textbf{zero perplexity degradation} ($\Delta\text{ppl} = 0.00$) with $80\%$ KV cache memory savings.
On compute-bound models ($7$B+), throughput retention reaches \textbf{98--101\%} of uncompressed baseline.
We also report a \textbf{principled negative result}: chunked attention via Python loops loses decisively to MLX's fused \texttt{mx.fast.scaled\_dot\_product\_attention} on Apple Silicon ($3$--$5\times$ slower at 16K--32K context), with bit-identical output. The implementation remains shipped for CUDA users without FlashAttention but should not be enabled on MLX. The total system spans 447 tests (zero regressions on the 293 pre-v0.4 tests), 12 of 13 referenced compression papers, and reproducible benchmarks for LLM throughput, video generation, and long-context memory.
\end{abstract}

% ============================================================================
% 1. INTRODUCTION
% ============================================================================
\section{Introduction}

KV cache compression is critical for enabling long-context inference on memory-constrained devices. The TurboQuant algorithm~\cite{zandieh2026turboquant} provides a theoretically optimal data-oblivious quantization scheme based on random orthogonal rotation followed by Lloyd-Max scalar quantization. However, the reference implementation leaves several systems engineering challenges unaddressed:

\begin{enumerate}[label=\arabic*.]
  \item \textbf{Reconstruction overhead}: Naive implementations dequantize the entire KV history on every token, creating $O(n^2)$ total work over $n$ tokens.
  \item \textbf{Python dispatch latency}: Multiple MLX/PyTorch operations per quantize/dequantize call incur interpreter overhead.
  \item \textbf{Codebook optimality}: Lloyd-Max codebooks are locally optimal but may not be globally optimal.
  \item \textbf{Diffusion transformer support}: TurboQuant was designed for autoregressive LLMs, not DiT architectures.
  \item \textbf{Architectural sprawl}: Each new compression paper from the literature (QuantSparse, DiTFastAttn, BSA, Fisher, Sheaf, etc.) requires modifying the same core files (\texttt{quantizer.py}, \texttt{cache/hf.py}, \texttt{cache/mlx.py}), making the library increasingly difficult to maintain and producing merge conflicts between contributors.
  \item \textbf{Production video pipeline support}: Diffusers video pipelines (WAN~2.2, LTX-2) crash or OOM on Apple Silicon for two unrelated reasons --- VAE decoder memory spikes and MPS \texttt{float64} incompatibility --- that block real local use even before any compression is considered.
\end{enumerate}

We address challenges 1--4 in \texttt{tqai} v0.3.1, challenge~5 with the v0.4 composable pipeline architecture, and challenge~6 with the v0.4 diffusers integration utilities. We also provide comprehensive benchmarks across 12 model configurations and report principled negative results for techniques that work in theory but lose to fused kernels on Apple Silicon.

% ============================================================================
% 2. BACKGROUND
% ============================================================================
\section{Background}

\subsection{PolarQuant / TurboQuant}

The TurboQuant Stage~1 (PolarQuant) quantization of a vector $\mathbf{x} \in \mathbb{R}^d$ proceeds as:

\begin{enumerate}[label=\arabic*.]
  \item \textbf{Norm extraction}: $s = \|\mathbf{x}\|_2$, stored as FP16.
  \item \textbf{Normalization}: $\tilde{\mathbf{x}} = \mathbf{x} / (s + \varepsilon)$.
  \item \textbf{Rotation}: $\mathbf{y} = \tilde{\mathbf{x}} \cdot R^\top$, where $R$ is a fixed Haar-distributed orthogonal matrix.
  \item \textbf{Scalar quantization}: For each coordinate~$i$, find
    \begin{equation}
      q_i = \operatorname*{argmin}_{j} \, |y_i - c_j|
    \end{equation}
    over Lloyd-Max centroids $\{c_1, \ldots, c_{2^b}\}$.
  \item \textbf{Storage}: $(q_1, \ldots, q_d)$ as uint8 indices $+\; s$ as float16.
\end{enumerate}

Dequantization reverses: lookup centroids $\to$ inverse rotate $\to$ scale by norm.

\begin{theorem}[Zandieh et al., 2026]
For $\mathbf{x} \sim \mathcal{N}(\mathbf{0}, I_d / d)$, the expected distortion of PolarQuant at $b$ bits per coordinate is:
\begin{equation}\label{eq:distortion}
  \mathbb{E}\bigl[\|\mathbf{x} - \hat{\mathbf{x}}\|^2\bigr] = D^*(b, d) + O(1/d)
\end{equation}
where $D^*(b, d)$ is the rate-distortion optimal distortion for the Gaussian source at rate~$b$.
\end{theorem}

\subsection{Lloyd-Max Quantization}

The Lloyd-Max algorithm finds optimal scalar quantization levels for a known distribution by iterating:

\begin{enumerate}[label=\arabic*.]
  \item \textbf{Centroid update}: $c_j = \mathbb{E}[X \mid X \in \text{partition } j]$.
  \item \textbf{Boundary update}: $b_j = (c_j + c_{j+1}) / 2$.
\end{enumerate}

For the post-rotation coordinate distribution $\mathcal{N}(0, 1/d)$, this converges to the minimum-MSE codebook.

\subsection{MLX and Apple Silicon}

MLX~\cite{apple2023mlx} is a machine learning framework for Apple Silicon featuring:
\begin{itemize}
  \item Unified memory architecture (no CPU$\leftrightarrow$GPU transfers).
  \item Lazy evaluation with automatic graph fusion.
  \item \texttt{mx.fast.metal\_kernel()} for custom Metal Shading Language (MSL) kernels.
\end{itemize}

% ============================================================================
% 3. INCREMENTAL CACHE RECONSTRUCTION
% ============================================================================
\section{Incremental Cache Reconstruction}

\subsection{Problem: $O(n^2)$ Full Reconstruction}

The naive cache implementation stores compressed entries and reconstructs the full history on every token:

\begin{algorithm}[H]
\caption{Naive Full Reconstruction}
\begin{algorithmic}[1]
\FOR{$t = 1, \ldots, n$}
  \STATE Store $\textsc{Compress}(\text{new\_kv})$
  \FOR{$i = 1, \ldots, t$}
    \STATE $\textsc{Decompress}(\text{entry}_i)$ \COMMENT{$O(t)$ work}
  \ENDFOR
\ENDFOR
\end{algorithmic}
\end{algorithm}

Total work:
\begin{equation}
  \sum_{t=1}^{n} t = \frac{n(n+1)}{2} = O(n^2)
\end{equation}

For $n = 100$ tokens at 24 layers $\times$ 2 (K+V), this is $100 \times 24 \times 2 \times 50.5 = 242{,}400$ dequantize calls.

\subsection{Solution: Incremental Buffer}

We maintain a running dequantized buffer alongside compressed storage:

\begin{algorithm}[H]
\caption{Incremental Buffer Reconstruction}
\begin{algorithmic}[1]
\FOR{$t = 1, \ldots, n$}
  \STATE $\text{entry} \gets \textsc{Compress}(\text{new\_kv})$
  \STATE $\text{buffer} \gets \textsc{Concat}(\text{buffer},\; \textsc{Decompress}(\text{entry}))$ \COMMENT{$O(1)$ dequant}
  \RETURN $\textsc{Concat}(\text{sink},\; \text{buffer})$
\ENDFOR
\end{algorithmic}
\end{algorithm}

Total work: $O(n)$ --- one dequantize per token.

\paragraph{Memory analysis.} Per KV vector:
\begin{itemize}
  \item Compressed: 130 bytes (128\,B uint8 indices $+$ 2\,B fp16 norm).
  \item Buffer: 256 bytes (fp16, $D=128$).
  \item Total: 386 bytes.
  \item vs.\ full reconstruction: 130\,B stored $+$ 512\,B fp32 materialized each call.
\end{itemize}

The incremental buffer uses \textbf{less} peak memory than full reconstruction.

\subsection{KIVI-style Residual Buffer}

Inspired by KIVI~\cite{yuan2024kivi}, we extend the incremental strategy with a residual window:

\begin{itemize}
  \item Last $R$ tokens kept uncompressed in full precision.
  \item Older tokens compressed $\to$ dequantized into incremental buffer.
  \item Assembly: $\textsc{Concat}(\text{sink},\; \text{old\_buffer},\; \text{recent\_uncompressed})$.
\end{itemize}

This provides zero quantization error on the $R$ most recent tokens, which receive the highest attention weights in most transformer architectures.

\subsection{Results}

\begin{table}[h]
\centering
\caption{Throughput improvement from incremental cache reconstruction (MLX, Apple Silicon).}
\label{tab:incremental-results}
\begin{tabular}{lccc}
\toprule
\textbf{Model} & \textbf{v0.2 (full)} & \textbf{v0.3.1 (incremental)} & \textbf{Speedup} \\
\midrule
Qwen 0.5B bf16 & 33 tok/s (10\%) & 118 tok/s (36\%) & $\mathbf{3.6\times}$ \\
Qwen 3B bf16   & 21 tok/s (30\%) & 66 tok/s (93\%)  & $\mathbf{3.1\times}$ \\
Llama 8B Q4    & 31 tok/s (30\%) & 85 tok/s (84\%)  & $\mathbf{2.7\times}$ \\
Qwen 7B Q8    & 26 tok/s (42\%) & 59 tok/s (93\%)  & $\mathbf{2.3\times}$ \\
Qwen 14B Q4   & 20 tok/s (35\%) & 50 tok/s (89\%)  & $\mathbf{2.5\times}$ \\
\bottomrule
\end{tabular}
\end{table}

% ============================================================================
% 4. FUSED METAL GPU KERNELS
% ============================================================================
\section{Fused Metal GPU Kernels}

\subsection{Kernel Design}

We fuse the entire quantize pipeline ($L_2$ norm, normalization, rotation, codebook search) into a single Metal compute kernel via \texttt{mx.fast.metal\_kernel()}.

\paragraph{Quantize kernel.} One thread per (vector, coordinate) pair. Grid $= (\text{num\_vecs} \times D,\; 1,\; 1)$.

Each thread:
\begin{enumerate}[label=\arabic*.]
  \item Computes the full $L_2$ norm redundantly ($D$ reads, $O(D)$ work).
  \item Normalizes its coordinate.
  \item Computes its rotated coordinate:
    \begin{equation}
      y^{(d)} = \sum_{j=0}^{D-1} \frac{x_j}{\|\mathbf{x}\|_2} \cdot R_{d,j}
    \end{equation}
  \item Performs serial argmin over $n_\text{levels}$ centroids.
  \item Stores uint8 index.
\end{enumerate}

The redundant norm computation trades $O(D)$ extra work per thread for avoiding shared memory (which \texttt{mx.fast.metal\_kernel} does not support). For $D \in \{64, 128\}$, the extra work is negligible compared to Python dispatch savings.

\paragraph{Dequantize kernel.} Reverse path --- centroid lookup, inverse rotation, norm scaling.

\subsection{Compile-Time Specialization}

We bake $D$ and $n_\text{levels}$ as compile-time constants into the MSL source:

\begin{lstlisting}[language=C]
const uint D = {D};
const uint N_LEVELS = {N_LEVELS};
\end{lstlisting}

This eliminates the need for \texttt{ensure\_row\_contiguous=False} (which injects per-call shape metadata), reducing kernel dispatch overhead.

\subsection{Correctness}

\begin{table}[h]
\centering
\caption{Metal kernel correctness: index match rate and maximum dequantization difference vs.\ Python path.}
\label{tab:kernel-correctness}
\begin{tabular}{cccc}
\toprule
\textbf{Bits} & \textbf{head\_dim} & \textbf{Index match} & \textbf{Max dequant diff} \\
\midrule
2 & 64, 128 & 100\%  & $1.2 \times 10^{-3}$ \\
4 & 64, 128 & 100\%  & $7.2 \times 10^{-7}$ \\
8 & 128     & 100\%  & $7.2 \times 10^{-7}$ \\
8 & 64      & 98.8\% & $1.2 \times 10^{-3}$ \\
\bottomrule
\end{tabular}
\end{table}

The 8-bit, $D=64$ case shows $<2\%$ index mismatch (off-by-1) due to FP rounding on closely-spaced centroids. This is benign --- adjacent centroids differ by ${\sim}0.001$.

% ============================================================================
% 5. EVOLUTIONARY CODEBOOK OPTIMIZATION
% ============================================================================
\section{Evolutionary Codebook Optimization}

\subsection{CMA-ES Solver}

We replace Lloyd-Max with Covariance Matrix Adaptation Evolution Strategy (CMA-ES; Hansen \& Ostermeier, 2001~\cite{hansen2001cmaes}), initialized from the Lloyd-Max solution.

\paragraph{Objective.} Minimize Monte Carlo estimate of expected distortion:
\begin{equation}\label{eq:cmaes-objective}
  \min_{\mathbf{c}} \; \hat{\mathbb{E}}\bigl[|X - Q_{\mathbf{c}}(X)|^2\bigr], \qquad X \sim \mathcal{N}(0, 1/d)
\end{equation}
where $Q_{\mathbf{c}}$ is the nearest-centroid quantizer with codebook $\mathbf{c} = (c_1, \ldots, c_{2^b})$.

\paragraph{Constraint.} Centroids sorted ascending after each generation (projection onto feasible set).

\paragraph{Results ($D=128$, 4-bit).}

\begin{table}[h]
\centering
\caption{Codebook solver comparison at $D=128$, 4-bit quantization.}
\label{tab:codebook-solvers}
\begin{tabular}{lccc}
\toprule
\textbf{Solver} & \textbf{MSE} & \textbf{Improvement} & \textbf{Time} \\
\midrule
Lloyd-Max           & $7.588 \times 10^{-5}$ & baseline      & 0.23\,s \\
Fuzzy C-means       & $7.585 \times 10^{-5}$ & $-0.04\%$     & 1.44\,s \\
CMA-ES              & $7.546 \times 10^{-5}$ & $\mathbf{-0.55\%}$ & 8.44\,s \\
\bottomrule
\end{tabular}
\end{table}

\subsection{Fuzzy C-Means with Temperature Annealing}

Inspired by Differentiable Soft Quantization~\cite{gong2019dsq}, we replace hard nearest-centroid assignment with temperature-controlled soft membership:
\begin{equation}\label{eq:fuzzy-membership}
  \mu_k(x) = \frac{\exp\!\bigl(-|x - c_k|^2 / \tau\bigr)}{\displaystyle\sum_{j} \exp\!\bigl(-|x - c_j|^2 / \tau\bigr)}
\end{equation}

Temperature annealing: $\tau_t = \tau_0 \cdot \alpha^t$, where $\alpha = 0.95$. At $\tau \to 0$, this converges to hard assignment.

\subsection{Attention-Aware Objective}

Inspired by APTQ~\cite{li2024aptq}, we provide an alternative objective that optimizes codebooks to preserve softmax attention distributions:
\begin{equation}\label{eq:attn-objective}
  \min_{\mathbf{c}} \; \mathbb{E}\!\left[\left\|\operatorname{softmax}\!\left(\frac{QK^\top}{\sqrt{d}}\right) - \operatorname{softmax}\!\left(\frac{Q\hat{K}^\top}{\sqrt{d}}\right)\right\|^2\right]
\end{equation}

This accounts for the fact that attention is invariant to uniform scaling and only cares about relative similarities.

% ============================================================================
% 6. CHUNKED ATTENTION VIA ONLINE SOFTMAX
% ============================================================================
\section{Chunked Attention via Online Softmax}

\subsection{Algorithm}

For sequence length $T_\text{kv}$ with chunk size $C$, we split $K, V$ into $\lceil T_\text{kv}/C \rceil$ chunks and maintain running accumulators using the online softmax algorithm~\cite{milakov2018online,dao2022flashattention}:

\begin{algorithm}[H]
\caption{Chunked Attention with Online Softmax}
\begin{algorithmic}[1]
\STATE \textbf{Initialize:} $m \gets -\infty$,\; $\ell \gets 0$,\; $O \gets \mathbf{0}$
\FOR{each chunk $c \in \{1, \ldots, \lceil T/C \rceil\}$}
  \STATE $S_c \gets Q \, K_c^\top \cdot \text{scale}$
  \STATE $m_\text{new} \gets \max(m,\; \max(S_c))$
  \STATE $\alpha \gets \exp(m - m_\text{new})$
  \STATE $P_c \gets \exp(S_c - m_\text{new})$
  \STATE $O \gets O \cdot \alpha + P_c \, V_c$
  \STATE $\ell \gets \ell \cdot \alpha + \sum P_c$
  \STATE $m \gets m_\text{new}$
\ENDFOR
\RETURN $O / \ell$
\end{algorithmic}
\end{algorithm}

\begin{theorem}
The output of~Algorithm~2 is identical to standard attention up to floating-point rounding.
\end{theorem}

\begin{proof}
Let $a_i$ denote the unnormalized attention weight for position~$i$, i.e., $a_i = \exp(\mathbf{q} \cdot \mathbf{k}_i / \sqrt{d})$. The standard output is
\begin{equation}
  O = \sum_{i} \frac{a_i}{\sum_{j} a_j} \cdot \mathbf{v}_i.
\end{equation}
The online algorithm maintains $\ell = \sum_{j} a_j \cdot \exp(m_\text{final} - m_j)$ where $m_j$ was the running max when $a_j$ was processed. Since $\exp(m_\text{final} - m_j) = \exp(m_\text{final})/\exp(m_j)$ and the corrections exactly cancel, $\ell = \sum_{j} a_j$. Similarly for $O$. \qedhere
\end{proof}

\subsection{Memory Analysis}

\begin{table}[h]
\centering
\caption{Attention matrix memory: full vs.\ chunked ($C=4096$).}
\label{tab:chunked-memory}
\begin{tabular}{lccc}
\toprule
\textbf{Tokens} & \textbf{Full attention} & \textbf{Chunked ($C\!=\!4096$)} & \textbf{Reduction} \\
\midrule
4,096     & 64\,MB  & 64\,MB (passthrough) & $1\times$ \\
48,360    & 17.8\,GB & 0.8\,GB & $\mathbf{22\times}$ \\
1,140,000 & 4.9\,TB  & 18\,GB  & $\mathbf{277\times}$ \\
\bottomrule
\end{tabular}
\end{table}

The last row represents WAN~2.2 at 720P $\times$ 81 frames through the VAE.

% ============================================================================
% 7. PALM: INFORMATION-THEORETIC ADAPTIVE BIT ALLOCATION
% ============================================================================
\section{Palm: Information-Theoretic Adaptive Bit Allocation}

\subsection{Novelty-Surprise Scoring}

For each token or activation tensor $\mathbf{x}$, we compute:

\paragraph{Novelty.}
\begin{equation}
  n(\mathbf{x}) = \frac{\|\mathbf{x} - \boldsymbol{\mu}_\text{EMA}\|}{\boldsymbol{\sigma}_\text{EMA}}
\end{equation}
where $\boldsymbol{\mu}_\text{EMA}$ and $\boldsymbol{\sigma}_\text{EMA}$ are exponential moving averages updated with decay $\alpha$.

\paragraph{Surprise.}
\begin{equation}
  s(\mathbf{x}) = \log\!\bigl(1 + \|\mathbf{x} - Q(\mathbf{x})\|^2\bigr)
\end{equation}
where $Q(\mathbf{x})$ is the quantized reconstruction.

\paragraph{Info score.}
\begin{equation}\label{eq:info-score}
  I(\mathbf{x}) =
  \begin{cases}
    \sqrt{n(\mathbf{x}) \cdot s(\mathbf{x})} & \text{if both available,} \\
    n(\mathbf{x}) & \text{otherwise.}
  \end{cases}
\end{equation}

\subsection{Tier Mapping}

\begin{table}[h]
\centering
\caption{Palm tier mapping from info score to bit width.}
\label{tab:palm-tiers}
\begin{tabular}{cccc}
\toprule
\textbf{Tier} & \textbf{Info score} & \textbf{Bits} & \textbf{Description} \\
\midrule
0 & $I < 0.1$           & 2 & Redundant --- aggressive compression \\
1 & $0.1 \leq I < 0.3$  & 3 & Expected --- standard compression \\
2 & $0.3 \leq I < 0.7$  & 4 & Novel --- preserve detail \\
3 & $I \geq 0.7$         & 8 & Critical --- minimal compression \\
\bottomrule
\end{tabular}
\end{table}

\subsection{Diffusion Step Scoring}

For diffusion transformers, the denoising schedule provides a natural info score:
\begin{equation}\label{eq:diffusion-score}
  I(t) = \frac{1}{1 + \text{SNR}(t)}
\end{equation}

Early steps (low SNR, high noise) $\Rightarrow$ $I \approx 1$ $\Rightarrow$ more bits.
Late steps (high SNR, refinement) $\Rightarrow$ $I \approx 0$ $\Rightarrow$ fewer bits.

This maps directly to the Lyapunov stability of the denoising trajectory: early steps are chaotic (sensitive to perturbation), late steps are stable (converging).

% ============================================================================
% 8. DIFFUSION TRANSFORMER INTEGRATION
% ============================================================================
\section{Diffusion Transformer Integration}

\subsection{Architecture Detection}

We extend the module detection system to recognize DiT attention patterns:

\begin{table}[h]
\centering
\caption{Attention projection attribute patterns across architectures.}
\label{tab:attn-patterns}
\begin{tabular}{lcccc}
\toprule
\textbf{Architecture} & \textbf{Q} & \textbf{K} & \textbf{V} & \textbf{Pattern} \\
\midrule
LLM (HuggingFace) & \texttt{q\_proj} & \texttt{k\_proj} & \texttt{v\_proj} & Standard \\
LLM (GPT-2)       & \texttt{c\_attn} (fused) & --- & --- & Fused \\
DiT (diffusers)   & \texttt{to\_q} & \texttt{to\_k} & \texttt{to\_v} & Diffusers \\
DiT (fused)       & \texttt{qkv} & --- & \texttt{proj} & Fused DiT \\
\bottomrule
\end{tabular}
\end{table}

\subsection{WAN 2.2 Integration}

WAN~2.2 TI2V-5B~\cite{wan2026} is a 5B-parameter text-image-to-video diffusion transformer. \texttt{tqai} detects 60 attention modules and 30 FFN modules in the WAN~2.2 transformer.

\paragraph{Text Encoder Cache.} The T5-XXL encoder output is fixed across all denoising steps. We cache it after the first call and return the cached version on subsequent calls, optionally compressed.

\paragraph{Inter-Step Delta Compression.} Adjacent denoising steps produce nearly identical K/V activations. We track the delta $\|K_t - K_{t-1}\| / \|K_t\|$ and store only the delta at low bits when it falls below a threshold.

\subsection{v0.4 Production Diffusers Integration}

The v0.3.x DiT integration achieved correct architecture detection but did not produce a usable end-to-end video generation experience on Apple Silicon. Two distinct blockers existed before any compression could be considered:

\paragraph{Blocker 1: VAE decoder memory spike.} WAN~2.2~5B's \texttt{AutoencoderKLWan.\_decode} accumulates 3D convolution feature maps in \texttt{self.\_feat\_map} (one entry per decoder conv layer) and a growing \texttt{out} tensor via \texttt{torch.cat} along the temporal axis. For 81~frames at $480 \times 832$, this allocates approximately 100~GB of intermediate tensors during the decode pass --- enough to OOM a 128~GB Mac. The diffusers VAE provides \texttt{enable\_tiling()} and \texttt{enable\_slicing()} methods, but they default to off and require explicit activation.

\paragraph{Blocker 2: MPS \texttt{float64} incompatibility (LTX-2).} The \texttt{LTX2RotaryPosEmbed1d} module in the diffusers LTX-2 implementation calls \texttt{torch.linspace(..., dtype=torch.float64, device=mps\_device)}, which crashes immediately because Apple's MPS backend does not support \texttt{float64}. The same issue exists in the LTX-2 transformer's RoPE module. The fix is a one-line attribute override (\texttt{module.double\_precision = False}) but it must be applied to two separate modules in the loaded pipeline.

v0.4 ships both fixes as drop-in utilities. After applying them, peak resident set size during a 33-frame WAN~2.2~5B generation is $32.2$~GB (model weights dominate), versus an unbounded spike that exceeds available memory at 81 frames without tiling.

\subsection{CFG Attention Sharing}

We implement DiTFastAttn's classifier-free guidance (CFG) sharing technique~\cite{yuan2024ditfastattn} as a strategy plugin (\texttt{cfg\_sharing}) that hooks into attention modules to cache the conditional pass output and serve it for the unconditional pass. We support both architectures observed in modern diffusers pipelines:

\begin{itemize}
  \item \textbf{Split-pass CFG (WAN~2.2):} Conditional and unconditional passes run as two separate forward calls. Hooks track phase via the pipeline's \texttt{cache\_context("cond"/"uncond")} and serve cached outputs during the unconditional phase. Achieves a $50\%$ share rate.
  \item \textbf{Batched CFG (LTX-2):} Both passes run in a single forward call with $[\text{uncond}, \text{cond}]$ concatenated along the batch dimension. Hooks copy the second half of the batch into the first half at each attention output. Achieves a $100\%$ share rate.
\end{itemize}

We measure the wall-clock impact in Section~\ref{sec:bench-video}.

\subsection{Forward Hook Limitation Resolved}

The v0.3.x device-transfer limitation (CPU$\leftrightarrow$MPS round-trips for the rotation matrix) is resolved in v0.4 by moving the rotation matrix to the input device on first use, plus shipping a parallel \texttt{MLXForwardCompressionHooks} class that uses MLX module replacement (\texttt{\_MLXCompressedWrapper}) instead of PyTorch's \texttt{register\_forward\_pre\_hook} API (which MLX does not provide). Forward compression on MLX now achieves the same correctness as PyTorch with no device-transfer overhead.

% ============================================================================
% 9. COMPOSABLE MIDDLEWARE PIPELINE (v0.4)
% ============================================================================
\section{Composable Middleware Pipeline (v0.4)}

\subsection{Motivation}

By v0.3.1, the library had absorbed a sequence of distinct techniques (incremental cache, residual cache, fused kernels, codebook solvers, forward hooks) by accreting flags onto \texttt{TurboQuantConfig} and modifying the same core files. Adding a new technique from the literature --- for example QuantSparse's second-order residual reparameterization~\cite{quantsparse2025} --- required modifying \texttt{quantizer.py} \emph{and} \texttt{cache/hf.py} \emph{and} \texttt{cache/mlx.py} simultaneously, with conflicts proportional to the number of in-flight contributors. The implementation also conflated two distinct concerns: \emph{what} to compress (per-token importance scoring) and \emph{how} to compress it (the actual quantization strategy), even though papers in the literature typically address only one or the other.

\subsection{Architecture}

We restructure the library around four protocol types defined in \texttt{pipeline/base.py}, each addressing one slot in the compression decision process: \texttt{Scorer} (per-token importance), \texttt{Strategy} (how to compress), \texttt{Monitor} (runtime parameter adjustment), and \texttt{ModelAdapter} (model family integration). Each protocol is a Python \texttt{Protocol} (PEP~544) with 3--4 methods. Plugins implement the protocol, register themselves by name in their directory's \texttt{\_\_init\_\_.py}, and become discoverable via \texttt{tqai plugins}. The runner (\texttt{pipeline/runner.py}) wraps the existing \texttt{PolarQuantizer} and dispatches \texttt{compress()} and \texttt{decompress()} through the optional scorer and strategy. When \texttt{pipeline=None} (the default), the runner short-circuits to direct \texttt{PolarQuantizer.quantize()} calls --- byte-identical to v0.3.1, zero overhead, all 293 pre-v0.4 tests pass byte-identically.

\subsection{Paper Playground}

We ship reference implementations of each technique from the recent compression literature as plugins. The implementations are intentionally self-contained and do not modify the proven v0.3.1 core. The current set covers 12 of 13 referenced papers:

\begin{table}[h]
\centering
\caption{Paper playground: each row is a single-file plugin.}
\label{tab:playground}
\small
\begin{tabular}{p{4cm}p{2.2cm}p{4cm}p{4cm}}
\toprule
\textbf{Paper} & \textbf{arXiv} & \textbf{Plugin} & \textbf{What it implements} \\
\midrule
TurboQuant (ICLR 2026) & 2504.19874 & \texttt{palm} + \texttt{tiered} & EMA novelty + dual-quantizer routing \\
Min-SNR (ICCV 2023) & 2303.09556 & \texttt{snr} scorer & Cosine + linear schedules \\
APTQ (DAC 2024) & 2402.14866 & \texttt{fisher} scorer & Squared-activation FIM proxy \\
Sheaf Attention (AAAI 2026) & 2601.21207 & \texttt{sheaf} scorer & Discrete Laplacian harmonicity \\
BSA (2025) & 2509.01085 & \texttt{bsa} scorer & Block centroid saliency (KV side) \\
DiTFastAttn step-share & 2406.08552 & \texttt{delta} strategy & First-order inter-step $\Delta$ \\
QuantSparse (2025) & 2509.23681 & \texttt{delta2} strategy & Second-order $\Delta^2$ with auto fallback \\
DiTFastAttn WA-RS & 2406.08552 & \texttt{window} strategy & Similarity-based output cache \\
DiTFastAttn CFG share & 2406.08552 & \texttt{cfg\_sharing} strategy & Split-pass + batched modes \\
Copresheaf NN (2025) & 2505.21251 & codebook registry & \texttt{head\_type} discriminator \\
Attention Analysis VDiTs & 2504.10317 & \texttt{skip\_layers} config & Per-layer compression bypass \\
SparseDiT (NeurIPS 2025) & 2412.06028 & tiered + skip\_layers & Tri-segment layer allocation \\
Spherical Attention & 2505.09326 & --- & Not implemented (requires retraining) \\
\bottomrule
\end{tabular}
\end{table}

\subsection{Plugin Authoring Cost}

A new paper requires (a)~one file in the appropriate directory, (b)~a one-line registration in the directory's \texttt{\_\_init\_\_.py}, and (c)~at least one test file. The proven core (\texttt{PolarQuantizer}, Lloyd-Max codebooks, cache strategies) is \emph{frozen}. Every measurement reported in Section~10 of the v0.3.1 paper revision ($\Delta\text{ppl}=0.00$ across 6 models) remains valid for the default path because that path is unchanged by construction.

\subsection{Genetic Algorithm Policy Search}

Because the pipeline configuration space is discrete (which scorer $\times$ which strategy $\times$ which monitor $\times$ per-strategy hyperparameters), we provide an offline GA search (\texttt{optimization/ga\_policy.py}) that evolves a \texttt{PolicyGenome} (9 genes) against an arbitrary fitness function. Population: 20. Generations: 10. Mutation rate: 0.15. Elitism: top $20\%$. Tournament selection: $k=3$. The user supplies an objective function (e.g., negative perplexity or NMSE) and the GA returns the best decoded pipeline configuration.

% ============================================================================
% 10. COMPREHENSIVE BENCHMARK RESULTS
% ============================================================================
\section{Comprehensive Benchmark Results}

\subsection{LLM Throughput (MLX, Apple Silicon, v0.4.0)}

\begin{table}[h]
\centering
\caption{LLM throughput with KV cache compression (MLX backend).}
\label{tab:llm-throughput}
\begin{tabular}{lcccccc}
\toprule
\textbf{Model} & \textbf{Params} & \textbf{Quant} & \textbf{Baseline} & \textbf{kv-only} & \textbf{Retention} & $\boldsymbol{\Delta}$\textbf{PPL} \\
\midrule
Qwen2.5-0.5B & 0.5B & bf16 & 302 tok/s & 88 tok/s  & 29\%  & $\mathbf{0.00}$ \\
Qwen2.5-3B   & 3B   & bf16 & 69 tok/s  & 24 tok/s  & 35\%  & $\mathbf{0.00}$ \\
Llama-3.1-8B & 8B   & Q4   & 28 tok/s  & 28 tok/s  & \textbf{100\%} & $\mathbf{0.00}$ \\
Qwen2.5-7B   & 7B   & Q8   & 17 tok/s  & 17 tok/s  & \textbf{101\%} & $\mathbf{0.00}$ \\
Qwen2.5-14B  & 14B  & Q4   & 15 tok/s  & 15 tok/s  & \textbf{98\%}  & $\mathbf{0.00}$ \\
Gemma 4 E4B  & 4B   & Q4   & 47 tok/s  & 24 tok/s  & 50\%  & $\mathbf{0.00}$ \\
\bottomrule
\end{tabular}
\end{table}

\textbf{Key finding:} Compute-bound models (7B+ with quantized weights) show \textbf{zero throughput cost} from KV cache compression. The incremental buffer overhead is completely absorbed by the model's own compute time.

\subsection{Pipeline Strategy Quality (synthetic NMSE)}

Synthetic streaming benchmark with 7 model profiles (Qwen 0.5B/3B/7B, Gemma 2B/7B, Llama~8B, WAN~2.2~5B), 5 inter-step time steps with decreasing perturbation noise (mimicking diffusion denoising), 64 sequence positions per step, 4 layers per model. Mean across all profiles:

\begin{table}[h]
\centering
\caption{v0.4 pipeline strategy quality (NMSE, lower is better).}
\label{tab:pipeline-nmse}
\begin{tabular}{lrrl}
\toprule
\textbf{Config} & \textbf{NMSE} & \textbf{vs Baseline} & \textbf{Notes} \\
\midrule
baseline                  & $0.009290$ & ---       & Direct PolarQuantizer at 4/2 bits \\
palm + tiered             & $0.009301$ & $+0.1\%$  & Adaptive bit allocation \\
\textbf{palm + delta}     & $\mathbf{0.003759}$ & $\mathbf{-59.5\%}$ & DiTFastAttn step-share \\
\textbf{snr + delta2}     & $\mathbf{0.003746}$ & $\mathbf{-59.7\%}$ & QuantSparse second-order $\Delta^2$ \\
sheaf + delta2            & $0.003763$ & $-59.5\%$ & Sheaf harmonicity + $\Delta^2$ \\
palm + window             & $0.018907$ & $+103.5\%$ & Cache reuse trades quality for speed \\
fisher + tiered           & $0.115858$ & $+1148\%$ & Proxy too aggressive --- offline mode required \\
skip\_layers              & $0.009297$ & $0.0\%$   & Layer protection passthrough \\
\bottomrule
\end{tabular}
\end{table}

\textbf{Key finding:} Inter-step delta strategies cut reconstruction distortion by ${\sim}60\%$ on synthetic streaming data. This empirically validates QuantSparse's claim~\cite{quantsparse2025} that the second-order residual is more compressible than the activation itself, and DiTFastAttn's claim~\cite{yuan2024ditfastattn} that adjacent denoising steps have high temporal redundancy.

\subsection{Video Generation: WAN 2.2 5B and LTX-2 (MLX)}\label{sec:bench-video}

End-to-end diffusers video generation, 33 frames at $480 \times 832$, 15 denoising steps, fixed seed:

\begin{table}[h]
\centering
\caption{Video generation benchmark (\texttt{tqai\_full} = \texttt{optimize\_vae\_memory} + \texttt{patch\_cfg\_sharing} + 8-bit forward hooks).}
\label{tab:video-benchmark}
\begin{tabular}{llrrrr}
\toprule
\textbf{Model} & \textbf{Config} & \textbf{Time} & \textbf{Peak RSS} & \textbf{CFG share} & \textbf{PSNR} \\
\midrule
WAN 2.2 5B & baseline   & 172.8\,s & 34.6 GB & $0\%$   & --- \\
WAN 2.2 5B & tqai\_full & 170.5\,s & 32.2 GB & $50\%$  & 56.5 dB \\
LTX-2      & baseline   & 88.3\,s  & 73.9 GB & $0\%$   & --- \\
LTX-2      & tqai\_full & 90.9\,s  & 74.4 GB & $100\%$ & 56.5 dB \\
\bottomrule
\end{tabular}
\end{table}

Two findings: (1)~\textbf{The actual local win is VAE memory optimization, not compression speedup.} Without \texttt{optimize\_vae\_memory()}, WAN~2.2~5B's VAE decoder spikes to ${>}100$~GB on an 81-frame video and OOMs on a 128~GB Mac. With it, peak resident set stays at 32~GB and 20-second videos fit comfortably. This single utility is the difference between ``video generation is impossible locally'' and ``video generation just works.''
(2)~\textbf{CFG sharing is functionally correct but does not deliver wall-clock speedup on MPS.} Hooks fire at the expected rate ($50\%$ on WAN's split-pass CFG, $100\%$ on LTX-2's batched CFG) and produce near-lossless output (PSNR 56.5~dB), but the Python hook dispatch overhead almost exactly cancels the saved attention compute on Apple Silicon's already-fast \texttt{mx.fast.scaled\_dot\_product\_attention}. We classify this as a \emph{functional but bottlenecked} result.

\subsection{Long-Context Chunked Attention (Honest Negative Result)}

Long-context generation benchmark on Qwen 2.5-3B (bf16) and Qwen 2.5-7B (8-bit), comparing baseline \texttt{mx.fast.scaled\_dot\_product\_attention} vs the chunked online-softmax implementation in \texttt{attention.py}:

\begin{table}[h]
\centering
\caption{Long-context chunked attention vs fused MLX SDPA.}
\label{tab:long-context}
\begin{tabular}{llrrrc}
\toprule
\textbf{Model} & \textbf{Context} & \textbf{Baseline} & \textbf{Chunked} & \textbf{Slowdown} & \textbf{Match} \\
\midrule
Qwen 3B bf16 & 4K  (no chunking) & 0.6\,s  & 0.6\,s  & $0\%$   & \checkmark \\
Qwen 3B bf16 & 8K                & 1.1\,s  & 3.3\,s  & $\mathbf{3.0\times}$ & \checkmark \\
Qwen 3B bf16 & 16K               & 2.6\,s  & 12.0\,s & $\mathbf{4.6\times}$ & \checkmark \\
Qwen 7B Q8   & 16K               & 5.8\,s  & 27.5\,s & $\mathbf{4.7\times}$ & \checkmark \\
Qwen 7B Q8   & 32K               & 17.4\,s & 68.3\,s & $\mathbf{3.9\times}$ & \checkmark \\
\bottomrule
\end{tabular}
\end{table}

\textbf{Output is bit-identical to baseline} on every configuration. The chunked attention math is correct. The slowdown is purely a Python dispatch problem: at 16K context with \texttt{chunk\_size=4096}, this is 4 chunks $\times$ ${\sim}10$ ops per chunk $\times$ 30 transformer layers $\times$ 2 attention modules per layer = approximately \textbf{2{,}400 separate Metal kernel launches per token}. The fused \texttt{mx.fast.scaled\_dot\_product\_attention} accomplishes the equivalent inner loop in \emph{one} kernel launch with hand-tuned tile sizes. Memory savings are also negligible at these scales: model weights (8.7~GB for Qwen 7B Q8) dominate the KV cache (well under 1~GB at 32K).

\textbf{We ship the chunked attention implementation but recommend against enabling it on MLX.} It is intended for CUDA users without FlashAttention, where the trade-off inverts.

\subsection{Codebook Quality}

\begin{table}[h]
\centering
\caption{Codebook solver comparison ($D=128$, 4-bit).}
\label{tab:codebook-quality}
\begin{tabular}{lcc}
\toprule
\textbf{Solver} & \textbf{MSE} & \textbf{Improvement} \\
\midrule
Lloyd-Max                  & $7.588 \times 10^{-5}$ & baseline \\
Fuzzy C-means ($\alpha=0.95$) & $7.585 \times 10^{-5}$ & $-0.04\%$ \\
CMA-ES (500 gen)           & $7.546 \times 10^{-5}$ & $\mathbf{-0.55\%}$ \\
\bottomrule
\end{tabular}
\end{table}

\subsection{Metal Kernel Correctness}

\begin{table}[h]
\centering
\caption{Metal kernel index match rate and maximum norm difference.}
\label{tab:metal-correctness}
\begin{tabular}{cccc}
\toprule
\textbf{Bits} & \textbf{head\_dim} & \textbf{Index match rate} & \textbf{Max norm diff} \\
\midrule
2 & 64, 128 & 100.0\% & 0.0 \\
4 & 64, 128 & 100.0\% & 0.0 \\
8 & 128     & 100.0\% & 0.0 \\
8 & 64      & 98.8\%  & 0.0 \\
\bottomrule
\end{tabular}
\end{table}

\subsection{Chunked Attention Accuracy}

All tests pass with $\text{atol} = 5 \times 10^{-4}$, $\text{rtol} = 2 \times 10^{-3}$ across:
\begin{itemize}
  \item Sequence lengths: 128, 256, 512, 1024.
  \item Chunk sizes: 32, 64, 128, 256, 512.
  \item Mask types: none, causal, additive.
  \item Head configurations: MHA ($H_q = H_{kv}$) and GQA ($H_q = 4 \cdot H_{kv}$).
\end{itemize}

% ============================================================================
% 10. TEST COVERAGE
% ============================================================================
\section{Test Coverage}

\begin{table}[h]
\centering
\caption{Test suite breakdown by component.}
\label{tab:test-coverage}
\begin{tabular}{lrl}
\toprule
\textbf{Component} & \textbf{Tests} & \textbf{Coverage} \\
\midrule
PolarQuantizer     & 46  & Round-trip, cosine, MSE, batch dims, zero vector \\
Metal kernels      & 26  & Bit-exact indices, norms, dequant, GQA \\
MLX cache strategies & 6 & Incremental, residual, full, sink tokens \\
HF cache strategies & 8  & All strategies + asymmetric bits \\
MLX forward hooks  & 7   & Attach/detach, shape/dtype, 4-bit vs 8-bit \\
Chunked attention  & 13  & Full vs chunked, causal, additive, GQA, chunk sizes \\
DiT detection      & 7   & \texttt{to\_q}/\texttt{to\_k}/\texttt{to\_v}, fused QKV, FeedForward, hook attachment \\
Palm scorer        & 6   & EMA tracker, tiers, bits, diffusion scoring \\
Step delta         & 3   & First step, delta usage, stats \\
Codebook solvers   & 6+  & Lloyd-Max, CMA-ES, fuzzy, symmetry, convergence \\
E2E large models   & 20+ & Real model inference with compression \\
\midrule
v0.4 Pipeline foundation     & 21 & base, registry, runner, builder, backward compat \\
v0.4 Scorers (palm, snr, fisher) & 23 & EMA, schedules, registration \\
v0.4 Strategies (tiered, delta, delta2, window) & 19 & Roundtrip, fallback chain, stats \\
v0.4 Sprint 4 (delta2, window, fisher, monitors) & 26 & Order fallback, similarity reuse, FTLE \\
v0.4 Adapters (llm, dit, wan) & 12 & Auto-detect, head info, registration \\
v0.4 Optimization (genome, GA) & 14 & Crossover, mutation, decode, fitness \\
v0.4 Paper gaps (sheaf, bsa, copresheaf, layer protection) & 22 & All new plugins \\
v0.4 CFG sharing & 8 & Split-pass + batched modes, cache lifecycle \\
\textbf{Total} & \textbf{447} & All 293 pre-v0.4 tests pass byte-identically \\
\bottomrule
\end{tabular}
\end{table}

% ============================================================================
% 11. RELATED WORK
% ============================================================================
\section{Related Work}

\begin{table}[h]
\centering
\caption{Comparison with related KV cache compression systems.}
\label{tab:related-work}
\begin{tabular}{llcc}
\toprule
\textbf{System} & \textbf{Approach} & \textbf{Throughput} & \textbf{Quality} \\
\midrule
TurboQuant~\cite{zandieh2026turboquant} & Rotation + Lloyd-Max & Not reported & Near-optimal \\
KIVI~\cite{yuan2024kivi}       & Asymmetric 2-bit, residual & $2.35$--$3.47\times$ & Minimal degrad.\ \\
KVQuant~\cite{hooper2024kvquant}    & LUT fused dequant-attn & $1.1$--$1.4\times$ & 1M ctx on A100 \\
GEAR~\cite{kang2024gear}        & Low-rank + sparse residual & $2.1$--$5.07\times$ & Near-lossless \\
BitDecoding~\cite{bitdecoding2025}  & Async pipeline, Tensor Core & Up to $8.6\times$ & $<3\%$ acc.\ drop \\
\textbf{tqai} (this work) & Incremental + Metal kernels & \textbf{98--101\% on 7B+} & $\boldsymbol{\Delta}\textbf{ppl} = \mathbf{0.00}$ \\
\bottomrule
\end{tabular}
\end{table}

% ============================================================================
% 12. CONCLUSION
% ============================================================================
\section{Conclusion}

\texttt{tqai} v0.4 makes three claims that we have validated empirically:

\begin{enumerate}[label=(\arabic*)]
  \item \textbf{The proven core is shippable.} K4/V2 KV quantization is byte-identical to baseline on every Qwen / Gemma / Llama model we tested ($\Delta\text{ppl} = 0.00$ across 6 models, 12 quantization variants). On compute-bound 7B+ models the throughput cost is zero. This has been true since v0.3.1; v0.4 does not change it by construction (\texttt{pipeline=None} short-circuits to v0.3.1 code paths, and all 293 pre-v0.4 tests pass byte-identically).
  \item \textbf{The plugin architecture lowers integration cost without raising baseline cost.} Adding a new compression paper from the literature now requires one new file in \texttt{scorers/}, \texttt{strategies/}, \texttt{monitors/}, or \texttt{adapters/}, plus a one-line registration. We demonstrate this by shipping reference implementations of 12 of 13 tracked papers without modifying \texttt{quantizer.py}, \texttt{cache/hf.py}, \texttt{cache/mlx.py}, or \texttt{kernels/}.
  \item \textbf{Practical local video generation requires solving non-compression problems first.} Two unrelated bugs in the diffusers stack (\texttt{AutoencoderKLWan} not enabling tiling by default, \texttt{LTX2RotaryPosEmbed1d} using \texttt{float64} on MPS) blocked any meaningful video work on Apple Silicon. We ship one-line fixes for both as drop-in utilities. After these utilities, WAN~2.2~5B generates 81-frame 480p video in 32~GB peak RSS instead of OOM-ing, and LTX-2 runs at all on Apple MPS instead of crashing immediately.
\end{enumerate}

We also report two \textbf{principled negative results} that we believe are worth more than the wins, because they save future contributors from dead ends:

\begin{itemize}
  \item \textbf{Chunked attention via Python loops loses to fused Metal kernels on MLX} by $3$--$5\times$ at 16K--32K context, with bit-identical output. The implementation is correct (15 tests prove it including a decode-mode regression test), the math is sound, and on CUDA without FlashAttention the speedup would be real. On Apple Silicon, \texttt{mx.fast.scaled\_dot\_product\_attention} already implements blocked attention internally and a Python-mediated loop cannot compete.
  \item \textbf{DiTFastAttn CFG sharing on MLX is functionally correct but does not deliver wall-clock speedup.} The hooks fire at the expected rate ($50\%$ / $100\%$), output PSNR is 56.5~dB (near-lossless), but the per-attention Python dispatch overhead almost exactly cancels the saved attention compute. The recommended next step is wrapping the hook in \texttt{mx.compile}; if that fails, a custom Metal kernel is required.
\end{itemize}

The unifying insight from both negative results: on Apple Silicon, \emph{anything that needs to fire inside the attention computation requires a fused kernel.} The pipeline cost of any new \texttt{scorer} or \texttt{strategy} that fires per-layer (not per-attention-call) is negligible because the framework itself is zero-overhead when unused and amortizes across the relatively cheap pre/post-attention path.

\subsection*{Future Work}

\begin{enumerate}[label=\arabic*.]
  \item \textbf{\texttt{mx.compile} the CFG sharing hook} --- would either close the speedup gap or deliver a definitive verdict on whether a custom Metal kernel is required.
  \item \textbf{Step distillation integration for video} --- distilled WAN/LTX checkpoints (LCM-style, Hyper-SD-style) reduce 20--50 denoising steps to 4--8 with no quality loss, delivering a $2.5$--$5\times$ video speedup that no compression technique on MLX has matched.
  \item \textbf{Long-context KV memory benchmark at 100K+ tokens} --- the K4/V2 win is currently demonstrated on quality, not memory; at 32K context the model weights still dominate the KV cache. At 100K+ context the KV cache becomes the primary memory consumer.
  \item \textbf{Offline gradient-based Fisher calibration mode} --- the runtime squared-activation proxy currently routes everything to the high-bit tier. A one-pass offline calibration would unblock the only broken plugin we ship.
  \item \textbf{Preset system} --- \texttt{tqai.patch(model, preset="...")} to remove the friction of picking scorer $\times$ strategy $\times$ monitor manually.
  \item \textbf{Custom Metal kernel for chunked SDPA at very long contexts ($>$64K)} --- only relevant if the fused \texttt{mx.fast.scaled\_dot\_product\_attention} itself OOMs at extreme sequence lengths.
  \item \textbf{On-device quantization for arbitrary dimensions} --- extend Metal kernels to non-power-of-2 head dimensions.
  \item \textbf{Spherical (L2-norm) attention} --- would require model retraining; out of scope for a drop-in compression library, but worth tracking as the only paper from Section~9.3 we did not ship.
\end{enumerate}

% ============================================================================
% REFERENCES
% ============================================================================
\bibliographystyle{unsrt}

\begin{thebibliography}{14}

\bibitem{zandieh2026turboquant}
A.~Zandieh, M.~Daliri, M.~Hadian, and V.~Mirrokni.
\newblock TurboQuant: Online Vector Quantization with Near-optimal Distortion Rate.
\newblock In \emph{ICLR}, 2026.
\newblock \href{https://arxiv.org/abs/2504.19874}{arXiv:2504.19874}.

\bibitem{yuan2024kivi}
J.~Yuan et~al.
\newblock KIVI: A Tuning-Free Asymmetric 2bit Quantization for KV Cache.
\newblock In \emph{ICML}, 2024.
\newblock \href{https://arxiv.org/abs/2402.02750}{arXiv:2402.02750}.

\bibitem{hooper2024kvquant}
C.~Hooper et~al.
\newblock KVQuant: Towards 10 Million Context Length LLM Inference with KV Cache Quantization.
\newblock In \emph{NeurIPS}, 2024.
\newblock \href{https://arxiv.org/abs/2401.18079}{arXiv:2401.18079}.

\bibitem{kang2024gear}
M.~Kang et~al.
\newblock GEAR: An Efficient KV Cache Compression Recipe for Near-Lossless Generative Inference of LLM.
\newblock In \emph{NeurIPS}, 2024.
\newblock \href{https://arxiv.org/abs/2403.05527}{arXiv:2403.05527}.

\bibitem{dao2022flashattention}
T.~Dao, D.~Y.~Fu, S.~Ermon, A.~Rudra, and C.~R\'{e}.
\newblock FlashAttention: Fast and Memory-Efficient Exact Attention with IO-Awareness.
\newblock In \emph{NeurIPS}, 2022.
\newblock \href{https://arxiv.org/abs/2205.14135}{arXiv:2205.14135}.

\bibitem{milakov2018online}
M.~Milakov and N.~Gimelshein.
\newblock Online normalizer calculation for softmax.
\newblock \href{https://arxiv.org/abs/1805.02867}{arXiv:1805.02867}, 2018.

\bibitem{hansen2001cmaes}
N.~Hansen and A.~Ostermeier.
\newblock Completely Derandomized Self-Adaptation in Evolution Strategies.
\newblock \emph{Evolutionary Computation}, 9(2), 2001.

\bibitem{gong2019dsq}
R.~Gong et~al.
\newblock Differentiable Soft Quantization: Bridging Full-Precision and Low-Bit Neural Networks.
\newblock \href{https://arxiv.org/abs/1908.05033}{arXiv:1908.05033}, 2019.

\bibitem{li2024aptq}
Y.~Li et~al.
\newblock APTQ: Attention-aware Post-Training Mixed-Precision Quantization for Large Language Models.
\newblock \href{https://arxiv.org/abs/2402.14866}{arXiv:2402.14866}, 2024.

\bibitem{apple2023mlx}
Apple.
\newblock MLX: An array framework for Apple silicon.
\newblock \url{https://github.com/ml-explore/mlx}, 2023.

\bibitem{wan2026}
Wan-AI.
\newblock WAN 2.2: Text-Image-to-Video Generation.
\newblock \url{https://huggingface.co/Wan-AI/Wan2.2-TI2V-5B}, 2026.

\bibitem{gemma4}
Google.
\newblock Gemma 4: Multimodal Language Models.
\newblock \url{https://huggingface.co/collections/google/gemma-4}, 2026.

\bibitem{idelgb2017}
Vector Quantization using the Improved Differential Evolution Algorithm for Image Compression.
\newblock \href{https://arxiv.org/abs/1710.05311}{arXiv:1710.05311}, 2017.

\bibitem{softquant2026}
Soft Quantization via Weight Coupling.
\newblock \href{https://arxiv.org/abs/2601.21219}{arXiv:2601.21219}, 2026.

\bibitem{bitdecoding2025}
BitDecoding.
\newblock \href{https://arxiv.org/abs/2503.18773}{arXiv:2503.18773}, 2025.

\bibitem{yuan2024ditfastattn}
Z.~Yuan, P.~Lu, H.~Zhang, X.~Ning, L.~Zhang, T.~Zhao, S.~Yan, G.~Dai, and Y.~Wang.
\newblock DiTFastAttn: Attention Compression for Diffusion Transformer Models.
\newblock In \emph{NeurIPS}, 2024.
\newblock \href{https://arxiv.org/abs/2406.08552}{arXiv:2406.08552}.

\bibitem{quantsparse2025}
QuantSparse: Second-Order Residual Quantization for Wan2.1 KV Cache.
\newblock \href{https://arxiv.org/abs/2509.23681}{arXiv:2509.23681}, 2025.

\bibitem{bsa2025}
BSA: Bidirectional Sparse Attention.
\newblock \href{https://arxiv.org/abs/2509.01085}{arXiv:2509.01085}, 2025.

\bibitem{minsnr2023}
T.~Hang et~al.
\newblock Efficient Diffusion Training via Min-SNR Weighting Strategy.
\newblock In \emph{ICCV}, 2023.
\newblock \href{https://arxiv.org/abs/2303.09556}{arXiv:2303.09556}.

\bibitem{cfg2022}
J.~Ho and T.~Salimans.
\newblock Classifier-Free Diffusion Guidance.
\newblock \href{https://arxiv.org/abs/2207.12598}{arXiv:2207.12598}, 2022.

\bibitem{vditanalysis2025}
Attention Analysis in Video Diffusion Transformers.
\newblock \href{https://arxiv.org/abs/2504.10317}{arXiv:2504.10317}, 2025.

\bibitem{sparsedit2025}
SparseDiT: Tri-Segment Token Allocation for Diffusion Transformers.
\newblock In \emph{NeurIPS}, 2025.
\newblock \href{https://arxiv.org/abs/2412.06028}{arXiv:2412.06028}.

\bibitem{sheaf2026}
Sheaf Attention \& Local Consistency.
\newblock In \emph{AAAI}, 2026.
\newblock \href{https://arxiv.org/abs/2601.21207}{arXiv:2601.21207}.

\bibitem{copresheaf2025}
Copresheaf Neural Networks.
\newblock \href{https://arxiv.org/abs/2505.21251}{arXiv:2505.21251}, 2025.

\bibitem{spherical2025}
Spherical Attention via Neural Circuit Diagrams.
\newblock \href{https://arxiv.org/abs/2505.09326}{arXiv:2505.09326}, 2025.

\bibitem{ltx2025}
Lightricks.
\newblock LTX-2: Long-Form Video Generation.
\newblock \url{https://huggingface.co/Lightricks/LTX-2}, 2025.

\end{thebibliography}

% ============================================================================
% APPENDICES
% ============================================================================
\appendix

\section{Mathematical Notation}\label{app:notation}

\begin{table}[h]
\centering
\begin{tabular}{cl}
\toprule
\textbf{Symbol} & \textbf{Meaning} \\
\midrule
$\mathbf{x} \in \mathbb{R}^d$        & Input KV vector \\
$R \in O(d)$                          & Haar-distributed orthogonal rotation matrix \\
$\mathbf{y} = \tilde{\mathbf{x}} \cdot R^\top$ & Rotated unit vector \\
$\{c_1, \ldots, c_L\}$               & Lloyd-Max centroids ($L = 2^b$) \\
$Q(\cdot)$                            & Quantization function (nearest centroid) \\
$s = \|\mathbf{x}\|_2$               & $L_2$ norm (stored as FP16) \\
$D^*(b, d)$                           & Rate-distortion optimal distortion \\
$\alpha$                              & EMA decay parameter (Palm scorer) \\
$\tau$                                & Temperature parameter (fuzzy C-means) \\
$C$                                   & Chunk size (chunked attention) \\
$R$ (KIVI)                            & Residual window size \\
\bottomrule
\end{tabular}
\end{table}

\section{Implementation Statistics (v0.4)}\label{app:stats}

\begin{table}[h]
\centering
\begin{tabular}{lrr}
\toprule
\textbf{Metric} & \textbf{v0.3.1} & \textbf{v0.4} \\
\midrule
Total source lines                & ${\sim}4{,}500$ & ${\sim}7{,}500$ \\
Test count                        & 293 & \textbf{447} \\
Pre-v0.4 tests passing byte-identically & --- & \textbf{293 / 293} \\
Models benchmarked (LLM)          & 6 arch., 12 var.\ & unchanged \\
Models benchmarked (DiT)          & WAN 2.2 (synth.) & WAN 2.2 + LTX-2 (real) \\
Supported frameworks              & PyTorch, MLX & unchanged \\
Supported LLM architectures       & Llama, Qwen, Gemma, GPT-2 & unchanged \\
Supported DiT architectures       & --- & WAN 2.2, LTX-2, generic diffusers \\
Metal kernel count                & 2 & unchanged \\
Codebook solvers                  & 3 & unchanged \\
Cache strategies                  & 3 & unchanged \\
\textbf{Pipeline scorers}         & 0 & \textbf{5} (palm, snr, fisher, sheaf, bsa) \\
\textbf{Pipeline strategies}      & 0 & \textbf{5} (tiered, delta, delta2, window, cfg\_sharing) \\
\textbf{Pipeline monitors}        & 0 & \textbf{2} (stability, lyapunov) \\
\textbf{Model adapters}           & 0 & \textbf{3} (llm, dit, wan) \\
\textbf{GA optimizer genes}       & 0 & \textbf{9} \\
\textbf{DiT utilities}            & 0 & \textbf{3} (optimize\_vae\_memory, patch\_mps\_compatibility, patch\_cfg\_sharing) \\
\textbf{Papers covered (of 13)}   & 0 & \textbf{12} \\
Benchmark scripts                 & 1 & 4 \\
Reports                           & 0 & 3 (markdown + HTML) \\
PyPI package size                 & ${\sim}80$\,KB & ${\sim}150$\,KB \\
\bottomrule
\end{tabular}
\end{table}

\end{document}
